% Curated from paper commit 7ee9157; original equations and protocol identities retained.
\section{Exact gradients and local learning rules}
\label{note:exact_rules}

\subsection{Exact gradients in a directed conductance tree}
\label{note:derivation}

To determine which factors a local learning rule can compute, we first derive
the exact conductance gradient, following the factorization and LocalCA rule introduced in our predecessor preprint \citep{safaai2026localcredit}. This requires the forward steady state, the
local synaptic sensitivity and transport of the somatic loss derivative
through the tree.

\subsection{Steady-state compartment voltage}

Consider a rooted dendritic tree with soma node 0. Let $V_n$ and $C_n$ be
the voltage and capacitance of compartment $n$, and let $\mathcal C_n$ be
its set of children. A non-somatic compartment \(n\) receives excitatory contacts \(\mathcal S_n^{\mathrm E}\) and inhibitory
contacts \(\mathcal S_n^{\mathrm I}\), with
\(\mathcal S_n=\mathcal S_n^{\mathrm E}\mathbin{\dot\cup}
\mathcal S_n^{\mathrm I}\). A contact of type
\(X\in\{\mathrm E,\mathrm I\}\) has nonnegative maximal conductance \(g_i^X\),
nonnegative presynaptic activity \(x_i^X\), open conductance \(g_i^Xx_i^X\),
and reversal \(E_X\). We use
\(E_i^{\mathrm{rev}}=E_{\mathrm E}\) on
\(\mathcal S_n^{\mathrm E}\) and
\(E_i^{\mathrm{rev}}=E_{\mathrm I}\) on
\(\mathcal S_n^{\mathrm I}\); when contact type is not at issue, we write the
generic \(g_i^{\mathrm{syn}}\), \(x_i\), and \(E_i^{\mathrm{rev}}\).
Define the aggregate open conductances
\[
G_n^{\mathrm E}=\sum_{i\in\mathcal S_n^{\mathrm E}}g_i^{\mathrm E}x_i^{\mathrm E},
\qquad
G_n^{\mathrm I}=\sum_{j\in\mathcal S_n^{\mathrm I}}g_j^{\mathrm I}x_j^{\mathrm I}.
\]
The compartment also receives child-to-parent coupling conductances
\(g_{c\to n}^{\mathrm{den}}\), and has leak \(g_n^{\mathrm L}\) with reversal
\(E_{\mathrm L}\). A child
transmits $a_c=f_c(V_c)$, where $f_c$ is the branch output nonlinearity
(named ``reactivation'' in the implementation) and is the identity when that
mechanism is disabled. Before taking the steady state, the directed
compartmental dynamics are
\begin{equation}
C_n\dot V_n
=g_n^{\mathrm L}(E_{\mathrm L}-V_n)
+G_n^{\mathrm E}(E_{\mathrm E}-V_n)
+G_n^{\mathrm I}(E_{\mathrm I}-V_n)
+\sum_{c\in\mathcal C_n}g_{c\to n}^{\mathrm{den}}(a_c-V_n).
\label{eq:s_current_balance}
\end{equation}
The total child coupling is
\(G_n^{\mathrm{child}}=\sum_{c\in\mathcal C_n}g_{c\to n}^{\mathrm{den}}\)
and its voltage-weighted contribution is
\(J_n^{\mathrm{child}}=\sum_{c\in\mathcal C_n}g_{c\to n}^{\mathrm{den}}a_c\).
Setting $\dot V_n=0$ and collecting the coefficients of $V_n$ then gives
\begin{align}
g_n^{\mathrm{tot}}
&=g_n^{\mathrm L}+G_n^{\mathrm E}+G_n^{\mathrm I}
  +G_n^{\mathrm{child}},\\
V_n
&=R_n^{\mathrm{tot}}
\left(g_n^{\mathrm L}E_{\mathrm L}
+G_n^{\mathrm E}E_{\mathrm E}
+G_n^{\mathrm I}E_{\mathrm I}
+J_n^{\mathrm{child}}\right),
\qquad R_n^{\mathrm{tot}}=(g_n^{\mathrm{tot}})^{-1}.
\label{eq:s_voltage}
\end{align}
Here $g_n^{\mathrm{tot}}$ and $R_n^{\mathrm{tot}}$ are the total
conductance and its reciprocal, respectively. The soma obeys the same current balance with $n=0$, including its child
couplings and any explicitly modeled somatic inputs, but has no parent term.
The controlled regular-tree architectures used here have no external somatic
input synapses unless a control is explicitly described otherwise.
For nonnegative conductances and activities, Eq.~\ref{eq:s_voltage} is a
conductance-weighted average of leak, synaptic reversal potentials and child
activities. This boundedness is useful for numerical stability, but is not
required for the chain-rule factorization below. Both E and I increase
\(g_n^{\mathrm{tot}}\); they differ through the driving forces
\(E_{\mathrm E}-V_n\) and \(E_{\mathrm I}-V_n\), not through a signed
conductance. When \(E_{\mathrm I}<V_n<E_{\mathrm E}\), the former is
depolarizing and the latter hyperpolarizing. At \(V_n=E_{\mathrm I}\), an
inhibitory conductance has zero net current but still reduces
\(R_n^{\mathrm{tot}}\), producing a pure shunt at that operating point.

The directed artificial-tree simulations use the dimensionless convention
\(g_n^{\mathrm L}=1\), \(E_{\mathrm L}=E_{\mathrm I}=0\), and
\(E_{\mathrm E}=1\), so inhibition is denominator-only. The reconstructed
focal-cable analysis uses \(E_{\mathrm E}=1\) and
\(E_{\mathrm I}=-0.2\), so its inhibitory perturbation is both
hyperpolarizing and shunting. Neither convention is a fit in millivolts.

\subsection{Local sensitivities and synaptic eligibility}

Writing the numerator of Eq.~\ref{eq:s_voltage} as
$b_n=g_n^{\mathrm{tot}}V_n$, quotient differentiation gives
\begin{align}
\frac{\partial V_n}{\partial g_i^{\mathrm{syn}}}
&=\frac{x_iE_i^{\mathrm{rev}} g_n^{\mathrm{tot}}-x_i b_n}
        {(g_n^{\mathrm{tot}})^2}
=x_iR_n^{\mathrm{tot}}(E_i^{\mathrm{rev}}-V_n),
\label{eq:s_syn_sensitivity}\\
\frac{\partial V_n}{\partial V_c}
&=f_c'(V_c)g_{c\to n}^{\mathrm{den}}R_n^{\mathrm{tot}},
\label{eq:s_transfer_sensitivity}\\
\frac{\partial V_n}{\partial g_{c\to n}^{\mathrm{den}}}
&=R_n^{\mathrm{tot}}(a_c-V_n).
\label{eq:s_coupling_sensitivity}
\end{align}
For a differentiable task loss $\mathcal L$, the exact synaptic
conductance gradient therefore factors as
\begin{equation}
\frac{\partial\mathcal L}{\partial g_i^{\mathrm{syn}}}
=\underbrace{x_iR_n^{\mathrm{tot}}(E_i^{\mathrm{rev}}-V_n)}_{\text{local eligibility}}
\underbrace{\frac{\partial\mathcal L}{\partial V_n}}_{\text{compartment error}}.
\label{eq:s_factorization}
\end{equation}
The eligibility contains quantities available at the synapse or its host
compartment. All remaining task dependence lies in the compartment error. In
particular, excitatory and inhibitory contacts use
\(x_i^{\mathrm E}R_n^{\mathrm{tot}}(E_{\mathrm E}-V_n)\) and
\(x_j^{\mathrm I}R_n^{\mathrm{tot}}(E_{\mathrm I}-V_n)\), respectively.
Near \(V_n=E_{\mathrm I}\), the inhibitory contact's own gradient can be small
even though its conductance still reduces \(R_n^{\mathrm{tot}}\) and changes
the sensitivity of other inputs.

\subsection{Transport of the somatic error}

Suppose the loss depends on the tree through its somatic output. Let $p(n)$
be the unique parent of non-somatic node $n$, and define the somatic voltage
error $\delta_0^V=\partial\mathcal L/\partial V_0$. The chain rule gives
\begin{equation}
\frac{\partial\mathcal L}{\partial V_n}
=\frac{\partial\mathcal L}{\partial V_{p(n)}}
 f_n'(V_n)R_{p(n)}^{\mathrm{tot}}g_{n\to p(n)}^{\mathrm{den}}.
\label{eq:s_recursion}
\end{equation}
Because the graph is a tree, iterating Eq.~\ref{eq:s_recursion} along the
unique path from $n$ to the soma yields
\begin{align}
\widetilde\alpha_n
&=\prod_{(i\to k)\in\operatorname{path}(n\to0)}
 f_i'(V_i)R_k^{\mathrm{tot}}g_{i\to k}^{\mathrm{den}},\\
\frac{\partial\mathcal L}{\partial V_n}
&=\widetilde\alpha_n\delta_0^V.
\label{eq:s_path_gain}
\end{align}
The effective path gain $\widetilde\alpha_n$ includes every local transfer
derivative. With identity reactivation, it reduces to the conductance-stage
gain $\alpha_n^{\mathrm{cond}}$ obtained by removing the $f_i'$ factors.
Equation~\ref{eq:s_path_gain} is an exact derivative, not a separately fitted
model parameter.

\subsection{Raw-parameter transforms}

Trainable raw parameters $\theta$ are transformed to nonnegative conductances
by a monotone map $g=\psi(\theta)$. The optimizer receives
\begin{equation}
\frac{\partial\mathcal L}{\partial\theta}
=\frac{\partial\mathcal L}{\partial g}\psi'(\theta).
\end{equation}
Where $\psi'(\theta)>0$, this transform preserves the sign of the
conductance gradient. Active synapse masks are applied after forming the
conductance-space gradient. The exact-gradient diagnostics compare gradients
in the same parameter space and include the transform derivative.

\subsection{General implicit-adjoint form}

The preceding path product is a corollary of a more general steady-state
identity. Let $\bm V$, $\bm g$ and $\bm x$ collect compartment voltages,
conductances and inputs. Suppose their steady state is defined by the
residual equations $\bm F(\bm V,\bm g,\bm x)=\bm 0$, with nonsingular
voltage Jacobian $J_V=\partial\bm F/\partial\bm V$. Define the adjoint
$\bm q$ by
\begin{equation}
J_V^{\mathsf T}\bm q=\nabla_{\bm V}\mathcal L.
\label{eq:s_general_adjoint}
\end{equation}
The implicit-function theorem gives
\begin{equation}
\frac{\partial\mathcal L}{\partial g_i}
=-\bm q^{\mathsf T}\frac{\partial\bm F}{\partial g_i}.
\label{eq:s_general_gradient}
\end{equation}
For a passive residual $\bm F=G\bm V-\bm b$, let $G$ be the conductance
matrix, $\bm b$ the vector of imposed current contributions and $\bm e_n$
the coordinate vector for compartment $n$. A conductance synapse at $n$ has
$\partial\bm F/\partial g_i=x_i(V_n-E_i^{\mathrm{rev}})\bm e_n$, and therefore
\begin{equation}
\frac{\partial\mathcal L}{\partial g_i}
=x_i(E_i^{\mathrm{rev}}-V_n)q_n.
\label{eq:s_general_synapse}
\end{equation}
Both contact classes add positive diagonal loading to \(G\), while their
different reversals enter \(\bm b\) and the local driving force in
Eq.~\ref{eq:s_general_synapse}.
Equations~\ref{eq:s_general_adjoint}--\ref{eq:s_general_synapse} recover the
directed path product when $J_V$ is triangular. For the reciprocal passive
cable, $J_V=G$ is symmetric positive definite, so
$G^{-\mathsf T}=G^{-1}$. A general nonsymmetric or recurrent operator instead
requires the inverse transpose. Active steady states enter through their
Jacobian whenever the required differentiability and nonsingularity conditions
hold; no passive symmetry is assumed by the general identity.

\subsection{Conductance-specific control of path gain}

Holding coupling conductances and the forward state fixed, add effective
inhibitory conductance $\Delta G_k^{\mathrm I}\geq0$ at compartment $k$. This
is an activity-weighted aggregate conductance, not an individual maximal
synaptic conductance. For a compartment $n$
whose path crosses $k$,
\begin{equation}
\frac{\alpha_n^{\mathrm{cond}}(\Delta G^{\mathrm I})}
     {\alpha_n^{\mathrm{cond}}(0)}
=\prod_{k\in\mathcal A(n)}
 \frac{g_k^{\mathrm{tot}}}{g_k^{\mathrm{tot}}+\Delta G_k^{\mathrm I}},
\qquad
\frac{\partial\log\alpha_n^{\mathrm{cond}}}{\partial G_k^{\mathrm I}}
=\begin{cases}
-R_k^{\mathrm{tot}}, & k\in\mathcal A(n),\\
0, & k\notin\mathcal A(n),
\end{cases}
\label{eq:s_inhibitory_gain}
\end{equation}
where
$\mathcal A(n)=\{k:(i\to k)\in\operatorname{path}(n\to0)\}$ contains the
parent endpoints of the path edges: the strict ancestors of $n$, including
the soma if its conductance is varied. Equation~\ref{eq:s_inhibitory_gain} proves path-selective
attenuation. It does not prove that attenuation improves optimization.
No reversal potential appears in Eq.~\ref{eq:s_inhibitory_gain} because this
calculation retains only the total-conductance denominator while holding the
forward state fixed. Any positive conductance has the same direct denominator
effect. It is specifically a shunt when \(E_{\mathrm I}\) is near the local
voltage, so that its accompanying current is small; the full forward and
gradient effects retain \(E_{\mathrm I}\).

Let $z_n=\log\alpha_n^{\mathrm{cond}}$ and let $\beta_n$ be the
accumulated log-gain attenuation along path $n$. The perturbed log gain is
$z_n'=z_n-\beta_n$. Across any specified ensemble of paths,
\begin{equation}
\Var(z')=\Var(z)+\Var(\beta)-2\Cov(z,\beta).
\label{eq:s_compression}
\end{equation}
Shunting narrows this distribution only if
$\Cov(z,\beta)>\Var(\beta)/2$. Path-gain compression is therefore an
operating-point-dependent empirical outcome rather than a theorem.

\subsection{Local rules and feedback objects}
\label{note:rules}

The exact compartment error derived above is generally unavailable to a local
synapse. We therefore retain the exact local eligibility and ask what
restricted teaching field can replace the transported error.

\subsection{Three-factor update}

The three-factor rule retains the exact local eligibility and substitutes an
available voltage-space teaching field $\widehat\delta_n$ for the exact
compartment error:
\begin{equation}
\widehat\nabla_{g_i^{\mathrm{syn}}}\mathcal L
=\left\langle x_iR_n^{\mathrm{tot}}(E_i^{\mathrm{rev}}-V_n)
\widehat\delta_n\right\rangle_B,
\qquad
\widehat\nabla_{g_{c\to n}^{\mathrm{den}}}\mathcal L
=\left\langle R_n^{\mathrm{tot}}(a_c-V_n)
\widehat\delta_n\right\rangle_B,
\label{eq:s_3f}
\end{equation}
where $B$ is the number of examples in a minibatch and
$\langle\cdot\rangle_B$ denotes their average. Excitatory and
inhibitory contacts use the same expression; their reversal potentials set the
sign and magnitude of the driving force.

The four- and five-factor rules multiply Eq.~\ref{eq:s_3f} by slowly varying,
bounded stage-level scalar reliability estimates. For stage $\ell$, mini-batch $t$, and example
$b$, the implementation averages the recorded layer voltage and corresponding
somatic population activity over their feature axes, giving $\bar V_{\ell,b}^{(t)}$ and
$\bar V_{0,b}^{(t)}$. The runs reported here used an uncentered dot product,
\begin{align}
s_{\ell,t}&=\frac{1}{B}\sum_{b=1}^{B}
 \bar V_{\ell,b}^{(t)}\bar V_{0,b}^{(t)},\\
m_{\ell,1}&=s_{\ell,1},\qquad
m_{\ell,t}=0.9m_{\ell,t-1}+0.1s_{\ell,t}\quad(t>1),\\
r_{\ell,t}^{\mathrm{4F}}&=\operatorname{clip}_{[0.1,2.0]}(m_{\ell,t}).
\label{eq:s_4f}
\end{align}
Thus $r_{\ell,t}^{\mathrm{4F}}$ is one scalar for a recorded stage, not a
separate branch-wise signal, and is not a centered covariance or a Pearson
correlation. Recorded voltages are detached from automatic differentiation,
and the same minibatch supplies $s_{\ell,t}$ and the local update. The positive
clamp prevents a negative or very small dot product from reversing or
eliminating the update.

The five-factor rule retains this multiplier and adds a conditional
variance-ratio factor. Let $Y_\ell^{(t)}$ be the recorded stage voltage and
$X_\ell^{(t)}$ its source-resolved parent proxy, each centered across mini-batch
$t$.
Let $\mathcal E_{0.1}[\cdot]$ denote an exponential moving average (EMA)
with update rate 0.1. Define
\begin{align}
v_{Y,\ell,t}&=\mathcal E_{0.1}[\langle (Y_\ell^{(t)})^2\rangle],&
v_{X,\ell,t}&=\mathcal E_{0.1}[\langle (X_\ell^{(t)})^2\rangle],&
c_{YX,\ell,t}&=\mathcal E_{0.1}[\langle Y_\ell^{(t)}X_\ell^{(t)}\rangle],\\
\beta_{\ell,t}&=\frac{c_{YX,\ell,t}}{v_{X,\ell,t}+10^{-3}},&
v_{\mathrm{res},\ell,t}&=\max\{v_{Y,\ell,t}-\beta_{\ell,t}c_{YX,\ell,t},10^{-12}\},\\
\phi_{\ell,t}&=\operatorname{clip}_{[0.25,4.0]}
 \left(\frac{v_{Y,\ell,t}}{v_{\mathrm{res},\ell,t}}\right),&
\widehat\nabla^{\mathrm{5F}}
&=r_{\ell,t}^{\mathrm{4F}}\phi_{\ell,t}\widehat\nabla^{\mathrm{3F}}.
\label{eq:s_5f}
\end{align}
If $v_{X,\ell,t}\leq10^{-12}$, the implementation sets $\phi_{\ell,t}=1$. These factors
are empirical preconditioners. They
are neither needed for the exact factorization nor extra task-error channels.
Under exact transported error, the three-factor rule recovers the exact
dendritic parameter gradient in the stated steady-state model. This separates
feedback approximation from the validity of the local eligibility term;
agreement of trained endpoints with backpropagation is tested empirically
rather than assumed from the identity.

\subsection{Feedback fields}

For neuron $u$ with $N_u$ routed dendritic compartments, indexed
$n=1,\ldots,N_u$ while the soma retains index $n=0$, let
$\bm\delta^V_u=(\delta^V_{1,u},\ldots,\delta^V_{N_u,u})^{\mathsf T}$ collect
the exact compartment errors and let $\widehat{\bm\delta}^V_u$ collect the
values available to the local rule. Let
$\bm c_u=(c_{u,1},\ldots,c_{u,K})^{\mathsf T}$ contain the signed values
carried by $K$ feedback channels. Let
$A_u\in\mathbb R^{N_u\times K}$ be their fixed compartmental address matrix:
column $k$ is the spatial field produced by unit amplitude on channel $k$, and
$(A_u)_{nk}$ is its delivery weight at compartment $n$. Linear superposition of
the channel fields gives
\begin{equation}
\widehat{\bm\delta}^{V}_u=A_u\bm c_u,
\qquad
\widehat\delta^V_{n,u}=\sum_{k=1}^{K}(A_u)_{nk}c_{u,k}.
\label{eq:s_routed_field}
\end{equation}
Equation~\ref{eq:s_routed_field} specifies the allowed feedback routes;
it is a modeling assumption separate from the membrane current balance. The address matrix specifies where signals can
be delivered; the coefficient vector specifies the values delivered on the
current example. The model is exact for $A_u=I$ and
$\bm c_u=\bm\delta^V_u$. It becomes a neuron-wide broadcast for
$K=1$ and $A_u=\bm 1_{N_u}$. A binary subtree column is one on descendants of
its branch point and zero elsewhere. The equation leaves coefficient
estimation open: projection analyses use best-fitting oracle coefficients to
measure representational capacity, whereas learning experiments use the
feedback rule specified for that task.

Here $K$ counts fixed spatial delivery profiles, not an intrinsic number of
independent external errors. Conditional on the forward state, a directed
single-output tree has the scalar somatic error multiplied by its
state-dependent path profile (Eq.~\ref{eq:s_path_gain}). Exact path transport
uses that scalar together with the sample-specific derivatives and is tested
as a separate transport resource. A context gate can likewise reside in the
local eligibility. The dimension of a field ensemble therefore measures the
demand on a fixed dictionary across states; it does not exclude scalar error
delivery combined with additional local transport or gating.

Table~\ref{tab:feedback_objects} defines the feedback objects without
conflating bandwidth with routing. The strict scalar field used for the
plotted MNIST ladder applies the scalar reduction defined below at every
local-rule stage. The matched-width scalar-fallback field
reuses a somatic vector only at a stage of the same width; at wider stages it
uses the same robust scalar reduction as the strict scalar condition. For an
$N$-dimensional somatic error vector $\bm\delta_b$ in example $b$, the code
first computes $\mu_b=N^{-1}\sum_u\delta_{b,u}$. With $\epsilon=10^{-6}$,
\begin{equation}
s_b=
\begin{cases}
\mu_b, & |\mu_b|>\epsilon,\\
\delta_{b,u^*},\quad u^*=\arg\max_u|\delta_{b,u}|, & |\mu_b|\leq\epsilon.
\end{cases}
\label{eq:s_scalar_reduction}
\end{equation}
For a one-dimensional error, the coordinate is retained. Selecting the
largest-magnitude coordinate when the mean nearly cancels prevents the
reduction from producing a near-zero teaching signal. The
replacement count was not recorded in the initial runs and was therefore
measured directly in the same-seed computational rerun. The matched-width field does not repeat each
neuron's coordinate over that neuron's descendants. The neuron-specific feedback field
does exactly that, using one communicated coordinate per neuron and the tree as
a fixed coordinate-to-arbor map. Exact path transport supplies the entire path derivative and
is an information upper bound.

In the regular $[3,3]$ networks, the matched-width scalar-fallback field retained the
128-dimensional somatic vector only at the 128-wide stage. It reduced the
field to one scalar at the 384- and 1,152-wide stages, so all 2,408,448
candidate feedforward synaptic parameters (92,160 active contacts) received
strict scalar feedback. The exception comprised 384 final branch-to-soma
couplings and 256 somatic gain and threshold parameters: 640 of 2,413,312 core
parameters (0.027\%). We therefore trained a separate 15-seed strict-scalar
condition in each architecture, which supplies the scalar rung in the original ladder retained in Supplementary Fig.~\ref{fig:si_mnist_dictionary_geometry}. The fresh six-rule ladder also uses strict scalar feedback.
Despite its small parameter count, this near-somatic exception affected
learning. Strict scalar reduced mean
accuracy relative to matched-width scalar-fallback feedback by 4.859 percentage points in shunting
networks (95\% paired-seed interval, 3.994--5.669; fallback higher in 15/15
seeds) and by 2.274 points in raw-additive networks (1.796--2.699; fallback higher
in 14/15 seeds, one tie). Both intervals lay well outside the prespecified
$\pm0.5$-point practical-equivalence margin. These results show that a small
near-somatic parameter group can exert disproportionate control over learning;
they do not imply that the remaining neurons collapse under strict scalar
feedback, because their local eligibility traces remain distinct.
The paired comparison is retained in the Source Data indexed in Table~\ref{tab:source_index}.

\subsection{Conditional gradient alignment}

For one tree and one example, suppose the correct somatic error is available
but path gain is omitted. Write local eligibility as $e_i$ and let $n(i)$ be
the host compartment. Then
\begin{equation}
g_i^{\mathrm{exact}}=e_i\widetilde\alpha_{n(i)}\delta_0^V,
\qquad
g_i^{\mathrm{shared}}=e_i\delta_0^V.
\end{equation}
Let $g^{\mathrm{exact}}$ and $g^{\mathrm{shared}}$ collect these synaptic
components, and let $\mathrm{CV}_w$ denote the weighted coefficient of
variation defined below. For nonnegative path gains and nonzero gradient vectors,
\begin{equation}
\cos(g^{\mathrm{shared}},g^{\mathrm{exact}})
=\frac{1}{\sqrt{1+\mathrm{CV}_w(\widetilde\alpha)^2}},
\qquad w_i=(e_i\delta_0^V)^2.
\label{eq:s_alignment}
\end{equation}
Here
\begin{equation}
\overline\alpha_w=\frac{\sum_iw_i\widetilde\alpha_{n(i)}}{\sum_iw_i},
\qquad
\mathrm{CV}_w(\widetilde\alpha)^2
=\frac{\sum_iw_i[\widetilde\alpha_{n(i)}-\overline\alpha_w]^2}
{\overline\alpha_w^2\sum_iw_i},
\end{equation}
so $\mathrm{CV}_w$ is the eligibility-energy-weighted coefficient of
variation of the path gain.
This identity explains why path-gain dispersion can matter after the correct
somatic coordinate has been supplied. It does not solve the separate problem
of supplying the correct coordinate across network layers.

\subsection{Topology-matched credit capacity}
\label{note:capacity_theory}

A restricted field is useful only if its spatial supports approximate the
required compartment-level credit. We therefore ask how much of that field can
be represented by $K$ ancestry-defined routes, and how this compares with a
point-neuron partition at the same bandwidth.

\subsection{Ancestry dictionary}

Let $A_T$ be the ancestry matrix of tree $T$, with one row per synaptic
coordinate and one column per candidate route. Let $\bm\beta$ contain the
state-dependent gains of these routes. A low-bandwidth teaching
coefficient vector $\bm c\in\mathbb R^K$ generates a synaptic field
\begin{equation}
\widehat{\bm g}=\Phi_{T,K}\bm c,
\qquad \Phi_{T,K}=D_eA_T\operatorname{diag}(\bm\beta)S_K,
\label{eq:s_dictionary}
\end{equation}
where the diagonal matrix $D_e$ holds local eligibility or anatomical weights
and $S_K$ assigns
the $K$ communicated channels to route sites. In this dictionary model,
topology fixes binary support and $\bm\beta$ rescales route columns. Nonzero
column rescaling preserves their span at a fixed selected route set, although
it can affect selection, conditioning and coefficient dynamics. The focal
cable interventions studied below are different: a physical shunt changes
spatial transfer and local eligibility, and need not be a diagonal gain on
these fixed columns.
The compartment-level matrix $A_u$ in Eq.~\ref{eq:s_routed_field} is the
general address map. Here $A_T$ is its anatomy-specific lift to synapse
coordinates, and $\Phi_{T,K}$ is the resulting synaptic-update dictionary
after eligibility, route gain and channel selection are included.

\subsection{Address dimension and the point-neuron resource comparison}

Separate the synapse-local eligibility from the task-dependent target credit
field and let $\bm t_u\in\mathbb R^{N_u}$ denote the latter for neuron $u$.
For this comparison, $\Phi$ acts on the chosen task-credit coordinates;
eligibility-weighted update dictionaries use the same projection construction
in their own coordinates. Let $W$ be a positive diagonal coordinate-weight
matrix, and let $\dagger$ denote the Moore--Penrose pseudoinverse. Define
\begin{equation}
\lVert\bm v\rVert_W^2=\bm v^{\mathsf T}W\bm v,
\qquad
P_{\Phi,W}=\Phi(\Phi^{\mathsf T}W\Phi)^\dagger\Phi^{\mathsf T}W.
\label{eq:s_weighted_projector}
\end{equation}
This is the $W$-orthogonal projector onto $\operatorname{col}(\Phi)$.
Equivalently, with $\overline\Phi=W^{1/2}\Phi$,
$\overline P_\Phi=\overline\Phi\,\overline\Phi^\dagger
=W^{1/2}P_{\Phi,W}W^{-1/2}$ is the ordinary Euclidean projector in the
transformed coordinates.
For a prespecified nested family of dictionaries $\Phi_{u,K}$, a residual
tolerance $0\leq\epsilon<1$ and positive expected target energy, define
\begin{equation}
d_{\epsilon,u}=\min\left\{K:
\frac{\mathbb E\lVert P_{\Phi_{u,K},W}\bm t_u\rVert_W^2}
{\mathbb E\lVert\bm t_u\rVert_W^2}\geq1-\epsilon\right\}.
\label{eq:s_address_dimension}
\end{equation}
We set $d_{\epsilon,u}=+\infty$ if no tested dictionary attains the threshold. This definition concerns the task-dependent coefficient after eligibility has
been factored out; heterogeneous presynaptic activity alone is therefore not
counted as an extra teaching coordinate. It is conditional on the model,
state ensemble, coordinate metric and dictionary family. It specifies a
resolution at a stated residual tolerance, rather than an intrinsic task
complexity or a unique optimal morphology.

For any fixed $\Phi\in\mathbb R^{N_u\times K}$, all fields generated as
$\Phi\bm c$ lie in $\operatorname{col}(\Phi)$ and hence occupy at most
$\operatorname{rank}(\Phi)\leq K$ dimensions. Conversely, the weighted
least-squares field in that space is $P_{\Phi,W}\bm t_u$, and its expected
weighted residual energy is
\begin{equation}
\mathbb E\lVert(I-P_{\Phi,W})\bm t_u\rVert_W^2
=\operatorname{tr}\!\left[(I-\overline P_\Phi)\overline\Sigma_t\right],
\qquad
\overline\Sigma_t=\mathbb E[W^{1/2}\bm t_u\bm t_u^{\mathsf T}W^{1/2}].
\label{eq:s_address_residual}
\end{equation}
Equations~\ref{eq:s_address_dimension}--\ref{eq:s_address_residual} prove the
fixed-dictionary dimension statement used in the main text;
$\overline\Sigma_t$ is an uncentered second moment. The ordinary point-neuron local
rule has $\Phi=\bm 1$ and therefore one task-dependent coefficient per neuron. A
$K$-route dendritic dictionary can use up to $K$ structured fields inside that
neuron. A point-neuron network is not mathematically forbidden from matching
them, but doing so requires $K$ separately routed teaching coefficients, $K$
gated parameter groups, or an equivalent expansion into additional units.
Accordingly, the comparison is about where address resources are represented,
not about an absolute expressive separation between all possible dendritic and
point-neuron networks.

\subsection{Field capture}

We now measure how much of a target field lies in a dictionary's span. Let
$\bm g$ denote that field in the coordinates under consideration, and use
the positive weighting from Eq.~\ref{eq:s_weighted_projector}. Set
$\overline{\bm g}=W^{1/2}\bm g$,
$\overline\Phi_{T,K}=W^{1/2}\Phi_{T,K}$, and let
$\overline P_\Phi$ be the ordinary Euclidean orthogonal projector onto the
columns of $\overline\Phi_{T,K}$. For a nonzero target field, capture is
\begin{equation}
\mathcal C_K(\bm g;\Phi,W)
=\frac{\lVert\overline P_\Phi\overline{\bm g}\rVert_2^2}
{\lVert\overline{\bm g}\rVert_2^2}
=1-\frac{\lVert\overline{\bm g}-\overline P_\Phi\overline{\bm g}\rVert_2^2}
{\lVert\overline{\bm g}\rVert_2^2}.
\label{eq:s_capture}
\end{equation}
The unweighted case has $W=I$. The MICrONS structural analysis uses summed
mapped excitatory synaptic area on the diagonal of $W$. For the transformed
uncentered second moment
$\overline\Sigma_g=\E[\overline{\bm g}\overline{\bm g}^{T}]$,
aggregate energy capture for a fixed projector is
\begin{equation}
\mathcal C_K(\overline\Sigma_g;\Phi,W)
=\frac{\operatorname{tr}(\overline P_\Phi\overline\Sigma_g)}
{\operatorname{tr}(\overline\Sigma_g)}.
\label{eq:s_covariance_matching}
\end{equation}
This is the ratio of expected projected energy to expected total energy. It
is generally different from the mean per-field fraction
$\E[\mathcal C_K(\bm g;\Phi,W)]$, which weights nonzero fields equally.
The two coincide for a fixed field norm, or when the second moment is formed
after normalizing each field. The MNIST dictionary gallery averages per-field
fractions; the cable-response operator analysis aggregates energy over its
columns. The second moment equals a centered covariance only for mean-zero
fields.
Equation~\ref{eq:s_covariance_matching} states the topology--task matching
condition: a sparse anatomical basis is useful only to the extent that task
credit lies in its span.

\begin{proposition}[Ky--Fan credit-capacity bound]
Let $\overline\Sigma_g\in\mathbb R^{p\times p}$ be positive semidefinite,
with eigenvalues satisfying
$\lambda_1\geq\cdots\geq\lambda_p\geq0$ and
$\operatorname{tr}(\overline\Sigma_g)>0$. For every rank-$K$ orthogonal
projector $P$,
\begin{equation}
\frac{\operatorname{tr}(P\overline\Sigma_g)}
{\operatorname{tr}(\overline\Sigma_g)}
\leq
\frac{\sum_{j=1}^{K}\lambda_j}{\sum_{j=1}^{p}\lambda_j}.
\label{eq:s_ky_fan_capture}
\end{equation}
Equality is attained by any rank-$K$ projector whose range is a
Ky--Fan-maximizing invariant subspace: it contains every eigenspace with
eigenvalue strictly greater than $\lambda_K$ and may choose the remaining
directions within the eigenspace at $\lambda_K$.
The difference between the right side and anatomical capture is the
rank-matched spectral regret.
\end{proposition}

\noindent\emph{Proof.}
Ky--Fan's maximum principle \citep{horn2013matrix} gives
$\max_{\substack{P=P^{\mathsf T}=P^2\\\operatorname{rank}(P)=K}}\operatorname{tr}(P\overline\Sigma_g)
=\sum_{j=1}^{K}\lambda_j$. Division by the fixed total spectral energy proves
Eq.~\ref{eq:s_ky_fan_capture}. \hfill$\square$

The spectral screen interpolates between two task-field second moments
$C_0$ and $C_1$ with equal trace, using
$C(\rho)=(1-\rho)C_0+\rho C_1$ for $0\leq\rho\leq1$. Thus, for two fixed projectors
$P$ and $Q$, the normalized capture difference has the sign of
\begin{equation}
\Delta(\rho)=\operatorname{tr}[(P-Q)C_0]
+\rho\operatorname{tr}[(P-Q)(C_1-C_0)],
\end{equation}
so an interior crossover is determined exactly by its two endpoints. This
affine result distinguishes alignment from rank: increasing $K$ raises the
oracle bound, whereas increasing $\rho$ determines whether the anatomical
subspace approaches it.

\subsection{One-step descent bound}

When the target field is a true loss gradient, capture also determines a
one-step descent guarantee in the corresponding metric. Let $\bm w$ be
the parameter vector and $\bm g=\nabla_{\bm w}\mathcal L$. To express
the weighted calculation as Euclidean descent, define
\begin{equation}
\overline{\bm w}=W^{-1/2}\bm w,
\qquad
\bm w=W^{1/2}\overline{\bm w},
\end{equation}
and define
$\overline{\mathcal L}(\overline{\bm w})=
\mathcal L(W^{1/2}\overline{\bm w})$. Then
$\nabla_{\overline{\bm w}}\overline{\mathcal L}=W^{1/2}\bm g=
\overline{\bm g}$. If $\overline{\mathcal L}$ has $L_W$-Lipschitz gradients
and $\eta$ is the step size, the descent lemma gives
\begin{equation}
\overline{\mathcal L}(\overline{\bm w}-\eta\overline P_\Phi\overline{\bm g})
\leq\overline{\mathcal L}(\overline{\bm w})
-\eta\overline{\bm g}^{T}\overline P_\Phi\overline{\bm g}
+\frac{L_W\eta^2}{2}
\lVert\overline P_\Phi\overline{\bm g}\rVert_2^2.
\end{equation}
Orthogonal projection implies
$\overline{\bm g}^{T}\overline P_\Phi\overline{\bm g}
=\lVert\overline P_\Phi\overline{\bm g}\rVert_2^2$, hence
\begin{equation}
\overline{\mathcal L}(\overline{\bm w}-\eta\overline P_\Phi\overline{\bm g})
\leq\overline{\mathcal L}(\overline{\bm w})
-\eta\left(1-\frac{L_W\eta}{2}\right)
\lVert\overline P_\Phi\overline{\bm g}\rVert_2^2.
\label{eq:s_descent}
\end{equation}
At $\eta=1/L_W$, the ratio between the projected- and full-gradient guarantees
is Eq.~\ref{eq:s_capture}. In the original coordinates, the projected step is
\begin{equation}
\bm w^+=\bm w-\eta W^{1/2}\overline P_\Phi W^{1/2}\bm g,
\end{equation}
which is a particular preconditioned update rather than an ordinary Euclidean
gradient step in $\bm w$. The implementation was checked on 10,000 random
positive-definite quadratic losses: there were no bound violations, minimum
numerical slack was $2.67\times10^{-7}$, and the capture identity agreed to
$1.89\times10^{-15}$. This numerical test checks the Euclidean projector algebra, which
also applies after an explicit coordinate transformation; it does not validate
synaptic-area weighting as a physical parameter metric. The result is a
one-step guarantee only when the tested field is an exact gradient in the
stated coordinates. It is not a convergence or generalization theorem and
does not show that a biological circuit can infer the oracle coefficients.

\subsection{Signal--noise theory of restricted credit routing}
\label{note:credit_phase_theory}

Capture alone favors increasing route resolution, but a broader operator also
admits more stochastic error. We next formalize this signal--noise trade-off
and ask when restricting the communicated credit improves expected progress.

\subsection{Expected descent from routed-update moments}

Let $\mathcal L:\mathbb R^p\rightarrow\mathbb R$ have $L$-Lipschitz gradient,
with $L>0$, on a neighborhood containing every update segment under consideration. At a fixed
parameter state, write $\bm g=\nabla\mathcal L(\bm w)$ and let $\bm d$ be a
routed update estimator with finite second moment. Define
$\bm m=\E[\bm d]$ and
$V=\E[(\bm d-\bm m)(\bm d-\bm m)^{\mathsf T}]$. The expectation is over
the stated sampling or feedback noise, conditional on $\bm w$. The update is
\begin{equation}
\bm w^+=\bm w-\eta\bm d.
\label{eq:s_operator_update}
\end{equation}

\begin{proposition}[Routed-update moment bound]
For Eq.~\ref{eq:s_operator_update},
\begin{equation}
\E\mathcal L(\bm w^+)\leq\mathcal L(\bm w)
-\eta\bm g^{\mathsf T}\bm m
+\frac{L\eta^2}{2}\left[
\lVert\bm m\rVert_2^2+\operatorname{tr}(V)
\right].
\label{eq:s_operator_bound}
\end{equation}
Write $a=\bm g^{\mathsf T}\bm m$ and
$b=\lVert\bm m\rVert_2^2+\operatorname{tr}(V)$. For $b>0$, minimizing
the bound over nonnegative steps gives $\eta^*=[a]_+/(Lb)$. Provided the
stated smoothness condition holds along the update segments at this step,
the certified decrease is
\begin{equation}
U=\frac{[\bm g^{\mathsf T}\bm m]_+^2}
{2L\{\lVert\bm m\rVert_2^2+\operatorname{tr}(V)\}},
\qquad [a]_+=\max(a,0).
\label{eq:s_operator_utility}
\end{equation}
For $b=0$, $\bm d=0$ almost surely and we set $U=0$. A zero score means
that this bound certifies no positive one-step decrease at the evaluated
state; it does not establish failure of every later update or learning rule.
\end{proposition}

\noindent\emph{Proof.}
Apply the descent lemma \citep{nesterov2004convex} to $-\eta\bm d$, take expectations and use
$\E\lVert\bm d\rVert_2^2=\lVert\bm m\rVert_2^2+\operatorname{tr}(V)$.
Minimizing the resulting scalar quadratic over $\eta\geq0$ gives the
claim. \hfill$\square$

A fixed linear routing operator gives a useful special case. If
$\bm d=M(\bm g+\bm\xi)$ with deterministic $M\in\mathbb R^{p\times p}$,
$\E\bm\xi=0$ and $\operatorname{Cov}(\bm\xi)=\Sigma$, then
$\bm m=M\bm g$ and $V=M\Sigma M^{\mathsf T}$. Equations
\ref{eq:s_operator_bound}--\ref{eq:s_operator_utility} then recover the
credit-operator expressions, denoted $U(M)$ below. The identity, orthogonal
projectors and fixed block-gain screens use this specialization.

The more general moment form is needed for sample-dependent eligibilities
and gates. In the context-gated task, changing an unselected distractor
leaves the exact parameter gradient unchanged but changes a shared-feedback
update through that distractor's eligibility. A fixed matrix of the exact
parameter gradient therefore cannot represent every routed update. The
output-error-to-update map remains linear conditional on the complete
forward state; that conditional map is a different object. The factorial
reanalysis measures routed means and second moments directly and does not
require a fixed map from exact parameter gradients to routed updates.

At fixed infinitesimal update norm, $-\bm g$ is the steepest Euclidean descent
direction by Cauchy--Schwarz. Consequently, no routed local operator provides a
uniform first-order improvement over exact full-batch backpropagation on the
same parameters and objective. Equation~\ref{eq:s_operator_utility} instead
compares restricted or noisy estimators and exposes when rejecting stochastic
energy compensates for discarded signal.

\begin{corollary}[Projected stochastic crossover]
Consider the identity-curvature quadratic
$\mathcal L(\bm w)=\frac12\lVert\bm w-\bm w^*\rVert_2^2$, initialize at
$\bm w=0$, set $\bm g=-\bm w^*$, use unit step size, and let $P$ be an
orthogonal projector. The expected post-update losses of the unprojected and
projected unbiased gradients are
\begin{align}
\E\mathcal L[-(\bm g+\bm\xi)]
&=\tfrac12\operatorname{tr}(\Sigma),\\
\E\mathcal L[-P(\bm g+\bm\xi)]
&=\tfrac12\left[
\lVert(I-P)\bm g\rVert_2^2+\operatorname{tr}(P\Sigma P)
\right].
\end{align}
The projected estimator is better precisely when
\begin{equation}
\operatorname{tr}[(I-P)\Sigma]
>\lVert(I-P)\bm g\rVert_2^2,
\label{eq:s_projection_crossover}
\end{equation}
that is, rejected noise exceeds discarded signal. A local estimator using the
same $P$ with additional independent, zero-mean coefficient noise cannot
achieve a lower expected loss than backpropagation followed by $P$.
\end{corollary}

For completeness, at a common step $0<\eta<2$ the unprojected-minus-projected
expected loss is
\begin{equation}
\frac{\eta}{2}\left\{
\eta\operatorname{tr}[(I-P)\Sigma]
-(2-\eta)\lVert(I-P)\bm g\rVert_2^2
\right\}.
\end{equation}
Hence projection is better exactly when
$\operatorname{tr}[(I-P)\Sigma]>
[(2-\eta)/\eta]\lVert(I-P)\bm g\rVert_2^2$; the unit-step result above is the
special case $\eta=1$. Smaller common steps therefore require proportionally
more rejected noise to compensate for the discarded signal.

\subsection{Interior optimum along the nested route ladder}
\label{note:interior_optimum}

The scalar, neuron-specific and correctly routed ancestry fields in the
controlled hierarchy analyses span nested subspaces. Their orthogonal
projectors form the ladder
$P_{\mathrm{scalar}}\subset P_{\mathrm{neuron}}\subset
P_{\mathrm{ancestry},K}\subset I$. Along such a ladder the capture of
Eq.~\ref{eq:s_capture} is monotone by construction, but the guaranteed
utility is not. For an orthogonal projector with retained signal energy
$S=\bm g^{\mathsf T}P\bm g=\lVert P\bm g\rVert_2^2$ and admitted noise
$N=\operatorname{tr}(P\Sigma P)$, the optimal-step utility of the
credit-operator bound is
\begin{equation}
U(S,N)=\frac{S^2}{2L\,(S+N)} .
\label{eq:s_projector_utility}
\end{equation}
Let one rung up the ladder add marginal signal $\Delta S\geq0$ and marginal
noise $\Delta N\geq0$. Substituting into Eq.~\ref{eq:s_projector_utility}
and clearing denominators gives the exact marginal condition
\begin{equation}
U(S+\Delta S,\,N+\Delta N)>U(S,N)
\iff
\Delta S\left[(S+N)(2S+\Delta S)-S^2\right]>S^2\,\Delta N .
\label{eq:s_marginal_condition}
\end{equation}
The bracket is positive, so for small marginal additions the condition
reduces to a threshold on the marginal signal-to-noise ratio,
$\Delta S/\Delta N>S/(S+2N)$: the utility keeps rising while each added
route scale contributes proportionally more signal than the running
threshold, and falls once the added scales admit mostly noise. When the
signal spectrum decays faster than the admitted noise across ladder scales,
$U$ can attain an interior maximum. On a finite tested ladder, the marginal
condition must become unfavorable before the last rung; spectral decay alone
does not guarantee that crossing or a unique optimum. The displayed marginal
condition assumes $S>0$; if $S=0$, a rung with positive retained signal has
positive utility whenever its total second moment is finite and nonzero.
Equations~\ref{eq:s_projector_utility} and
\ref{eq:s_marginal_condition}, together with the interior-maximum
construction, are verified numerically in the released test suite
(\texttt{tests/test\_laminar\_optimality.py}).
Random, deranged and independently gained route families elsewhere in the
paper need not lie on this same nested ladder.

Two experiments evaluate this optimum empirically
(Table~\ref{tab:source_index}; analysis code is included
in the release). In the hierarchy-depth
routing experiment the initialization utility of the aligned tree, computed
per seed and task depth, placed its argmax at the trained loss minimum in
135 of 200 seed--task pairs, and the group-mean argmax matched the trained
optimum exactly at $H_{\mathrm c}=1,2,3$; at the full-rank boundary
$H_{\mathrm c}=4$ the one-step
bound was conservative and favoured $D_{\mathrm{r}}=2$, reflecting its single-step
horizon.

The trained subtree factorial provides a second, retrospective test. Its
full-batch moment score
uses the global logistic curvature bound and the same four-control comparator
as the accuracy contrast: deranged, depth-interleaved, random-sparse and
random-rank routes. The mean ancestry-minus-best-control utility was positive
at $K=4$, negative at $K=1,2$ and zero at $K=8$. Its contrast-maximizing
budget was $K=4$ in all 20 seeds, matching the observed accuracy-contrast
maximum in 13 of 20. Mean accuracy-contrast selection regret was 1.94
percentage points. Adding the hypothetical minibatch variance made the mean
$K=4$ utility contrast negative, reduced optimum agreement to 8 of 20 and
gave mean regret 2.75 points. Thus the score matched to the deterministic
training recipe recovers the mean intermediate-bandwidth ancestry peak,
while its seed-level selection remains imperfect. This retrospective result
is distinct from the raw within-family performance optimum at $K=8$ and
from prospective morphology selection. The comparator definitions, scores
and prediction errors are supplied in
\texttt{source\_data/review\_evidence\_reanalysis/}.

\subsection{Laminar route families and hierarchical task second moments}
\label{note:laminar_optimality}

The distinction between a point neuron and a dendritic tree can be stated
as a constraint on the feasible credit operators at fixed rank. A point
neuron addresses its synapses with the single shared column $\bm 1$: a
rank-one operator with uniform support. An ancestry dictionary uses structured columns whose supports form a
\emph{laminar} family: any two
supports are nested or disjoint. An unconstrained rank-$K$ feedback matrix
is the comparison class. For the balanced binary tree used in the
factorial, the span of the $K$-route ancestry dictionary equals the span of
the $K$ coarsest vectors of the orthonormal tree-Haar basis
$U=[\bm h_1,\dots,\bm h_n]$ (constant, then sibling-subtree differences,
coarse to fine) for every admissible budget $K\in\{1,2,4,\dots,n\}$.
Here $n$ is the number of leaf coordinates. Let $P_{A,K}$ denote the
orthogonal projector onto the selected ancestry span.

\begin{proposition}[Laminar capture of a hierarchical second moment]
Let the task-credit second moment be $C=U\Lambda U^{\mathsf T}$ with
$\Lambda=\operatorname{diag}(\lambda_1,\dots,\lambda_n)$ ordered so that
$\lambda_1\geq\lambda_2\geq\cdots\geq\lambda_n\geq0$ along the
coarse-to-fine Haar order, with $\operatorname{tr}(C)>0$. Then for every admissible budget $K$ the
ancestry dictionary attains the Ky--Fan bound,
\begin{equation}
\frac{\operatorname{tr}(P_{A,K}\,C)}{\operatorname{tr}(C)}
=\frac{\sum_{i\leq K}\lambda_i}{\operatorname{tr}(C)}
=\max_{\substack{P=P^{\mathsf T}=P^2\\\operatorname{rank}(P)=K}}\frac{\operatorname{tr}(P\,C)}{\operatorname{tr}(C)} .
\end{equation}
For a general second moment $C$ with eigenpairs $(\lambda_i,\bm v_i)$ ordered
by decreasing eigenvalue, the shortfall from the Ky--Fan bound decomposes
exactly as
\begin{equation}
\sum_{i\leq K}\lambda_i-\operatorname{tr}(P_{A,K}C)
=\sum_{i\leq K}\lambda_i\lVert(I-P_{A,K})\bm v_i\rVert_2^2
-\sum_{j>K}\lambda_j\lVert P_{A,K}\bm v_j\rVert_2^2 ,
\end{equation}
the leading-eigenspace energy outside the laminar span minus the trailing
energy that the span retains. It is therefore bounded above by the former
term. The shortfall is zero exactly when $P_{A,K}$ is Ky--Fan maximizing;
when $K<n$ and $\lambda_K>\lambda_{K+1}$, this is equivalent to the ancestry span
being the unique leading-$K$ eigenspace.
\end{proposition}

\begin{proof}
The ancestry span equals the span of the $K$ coarsest Haar vectors, which
under the stated ordering are the top-$K$ eigenvectors of $C$; a rank-$K$
orthogonal projector maximizes $\operatorname{tr}(PC)$ exactly on a top-$K$ eigenspace
(Ky--Fan). For general $C$, expanding
$\operatorname{tr}(P_{A,K}C)=\sum_j\lambda_j\lVert P_{A,K}\bm v_j\rVert_2^2$
and $\lVert P_{A,K}\bm v_i\rVert_2^2=1-\lVert(I-P_{A,K})\bm v_i\rVert_2^2$
for $i\leq K$ gives the stated decomposition. Ky--Fan's equality condition,
including arbitrary choices inside a degenerate cutoff eigenspace, gives the
last statement.
\end{proof}

This ordered hierarchical second moment is one sufficient class for which
ancestry attains the rank-$K$ spectral optimum. At a fixed $K$, the equality
condition is the leading-eigenspace condition in the proposition; it does not
require the entire second moment to have a Haar eigenbasis and does not
identify a unique physical tree. Hierarchical organization without the stated
coarse-to-fine spectral ordering is not sufficient. The
spectral-alignment experiment realizes both regimes: at full alignment and
$K=4$ ancestry attained the rank-matched Ky--Fan bound to numerical
precision, whereas at random alignment its advantage over sampled random
routes vanished (main text). The attainment, the shortfall identity and
the failure case with a fine-scale-dominant second moment are verified
numerically in \texttt{tests/test\_laminar\_optimality.py}.

\subsection{Point-neuron partition residual and route-coefficient error}

Let $t_i$ be the exact task-dependent target remaining after local
eligibility has been factored out, let $w_i>0$ be coordinate weights, and let
$\pi(i)$ assign synapse $i$ to a group that must share one teaching
coefficient. The best allowed field minimizes
$\sum_iw_i[t_i-c_{\pi(i)}]^2$.

\begin{proposition}[Shared-coefficient residual]
For every group $G$, the optimum is its weighted mean
\begin{equation}
c_G^*=\frac{\sum_{i\in G}w_it_i}{\sum_{i\in G}w_i},
\end{equation}
and the irreducible residual is
\begin{equation}
\mathcal R(\pi)=\sum_G\sum_{i\in G}w_i(t_i-c_G^*)^2.
\label{eq:s_partition_variance}
\end{equation}
Refining a partition cannot increase $\mathcal R$; it decreases the residual
only when an old group contains coefficient variation resolved by the new
groups.
\end{proposition}

\noindent\emph{Proof.}
Differentiate the separable weighted least-squares objective with respect to
each $c_G$. The refinement statement follows because the old solution remains
feasible when two or more refined coefficients are constrained to be equal.
\hfill$\square$

The ordinary point-neuron rule has one group per neuron. Dendritic ancestry
offers a particular nested sequence of refinements, but is helpful only if the
task coefficient varies across those subtrees. An explicitly grouped point
model can implement the same partition, which is why the resource-matched
point controls in the main article coincide with the dendritic implementation.

For a general dictionary $\Phi$ and positive diagonal $W$, let
$\bm c^*=\arg\min_{\bm c}\lVert\bm t-\Phi\bm c\rVert_W^2$ and let
$\widehat{\bm c}$ be the coefficients actually supplied by a local circuit.
Weighted projection orthogonality yields the exact decomposition
\begin{equation}
\lVert\bm t-\Phi\widehat{\bm c}\rVert_W^2
=\underbrace{\lVert\bm t-\Phi\bm c^*\rVert_W^2}_{\text{address / representation error}}
+\underbrace{\lVert \Phi(\widehat{\bm c}-\bm c^*)\rVert_W^2}_{\text{coefficient-estimation error}}.
\label{eq:s_field_error_decomposition}
\end{equation}
If the coefficient error is unbiased with covariance $\Omega_c$, the expected
second term is
$\operatorname{tr}(\Phi^{\mathsf T}W\Phi\,\Omega_c)$. The Gram matrix
$G_\Phi=\Phi^{\mathsf T}W\Phi$ therefore controls how coefficient error is expressed.
Estimating $\bm c^*$ is sensitive to noise along directions with small
positive Gram eigenvalues. Correlated route columns can likewise make
individual coefficients poorly identifiable even when their combined field
is stable. Address rank alone is
not a complete learning metric.

\subsection{Static route gains preserve span but can change conditioning}

\begin{proposition}[Route-wise static-gain invariance]
Let $D_\beta$ be an invertible diagonal matrix of fixed nonzero route gains. Then
\begin{equation}
\operatorname{col}(\Phi D_\beta)=\operatorname{col}(\Phi),\qquad
P_{\Phi D_\beta,W}=P_{\Phi,W}.
\label{eq:s_static_gain_span}
\end{equation}
Thus route-wise static scalar gains cannot change address rank, projection
residual or spectral capture. They can change the Gram matrix from
$\Phi^{\mathsf T}W\Phi$ to
$D_\beta^{\mathsf T}\Phi^{\mathsf T}W\Phi D_\beta$ and hence alter
conditioning and coefficient dynamics.
\end{proposition}

This proposition concerns nonzero scalar rescaling of existing route columns.
Spatially varying row gains, state-dependent gating and gains that reach zero
can change the effective span or active support; those cases are not covered
by Eq.~\ref{eq:s_static_gain_span}. Physical shunts also change the cable
Green's function. For a somatic adjoint source on a passive tree, that change
has an exact diagonal-gain representation in the baseline-weighted ancestry
partition associated with the shunt site
(Corollary~\ref{cor:shunt_ancestry_gain}). Its partition indicators are
differences of nested subtree indicators. The diagonal form does not extend
to an arbitrary selected overlapping dictionary, and the local synaptic
driving force can still vary within each block.

\subsection{Finite-time coefficient learning has a bias--variance crossover}

Equal span need not give equal learning trajectories. Let
$\bm f_t=\Phi\bm c_t$ be the field produced by coefficients $\bm c_t$
after update $t$, and let $\lambda_{\max}>0$ be the largest eigenvalue of
$\Phi\Phi^{\mathsf T}$. Normalizing ordinary coefficient descent by
this eigenvalue gives $H_\Phi=\Phi\Phi^{\mathsf T}/\lambda_{\max}$.
For normalized step size $\eta$ and noisy rate-based target estimates
$\bm g+\bm\xi_t$, with
$\E\bm\xi_t=0$ and
$\operatorname{Cov}(\bm\xi_t)=\sigma^2I/n$, where $n$ is the effective
sample size and $\sigma^2$ the single-observation noise scale, the field
iteration is
\begin{equation}
\bm f_{t+1}=\bm f_t-\eta H_\Phi(\bm f_t-\bm g-\bm\xi_t).
\label{eq:s_coefficient_field_iteration}
\end{equation}

\begin{proposition}[Exact finite-time coefficient risk]
Assume $\bm f_0=0$, $\bm g\in\operatorname{col}(\Phi)$ and let
$(\mu_j,\bm u_j)$ be the positive eigenpairs of $H_\Phi$, with
$\theta_j=\bm u_j^{\mathsf T}\bm g$. For independent isotropic noise and a stable step satisfying
$|1-\eta\mu_j|<1$ for every positive mode, the expected population loss
after $T$ updates is
\begin{equation}
\E\frac{\lVert\bm f_T-\bm g\rVert_2^2}{2}
=\frac12\sum_j\left[
(1-\eta\mu_j)^{2T}\theta_j^2
+\frac{\sigma^2}{n}
\frac{(\eta\mu_j)^2[1-(1-\eta\mu_j)^{2T}]}
{1-(1-\eta\mu_j)^2}
\right].
\label{eq:s_coefficient_finite_risk}
\end{equation}
For Gram-preconditioned descent, $H_\Phi=P_\Phi$ and all positive $\mu_j=1$;
therefore every dictionary with the same projector has the same field
trajectory under common observations.
\end{proposition}

\noindent\emph{Proof.}
Project Eq.~\ref{eq:s_coefficient_field_iteration} onto each eigenvector. The
deterministic error contracts by $1-\eta\mu_j$ per step. Independent innovations
have variance $(\eta\mu_j)^2\sigma^2/n$ and accumulate as a finite geometric
series. Summing half the modal second moments gives
Eq.~\ref{eq:s_coefficient_finite_risk}. Gram preconditioning replaces
$\Phi^{\mathsf T}\Phi$ by its pseudoinverse on the positive subspace, so the induced
field operator is $P_\Phi$. \hfill$\square$

Equation~\ref{eq:s_coefficient_finite_risk} makes the low-data boundary
explicit. Small $\mu_j$ retain more finite-time bias but admit less estimator
variance. For two operators with the same span, let $\Delta B$ be the difference in
the bias term and $\Delta V$ the difference in the coefficient of the
noise term. Their risk difference is $\Delta B+(\sigma^2/n)\Delta V$. If $\Delta B>0$ and $\Delta V<0$, the
ordering reverses at
$n^*=-\sigma^2\Delta V/\Delta B$: slow coordinates can help below $n^*$ by
implicit noise filtering, whereas well-conditioned coordinates help above it
by reducing bias. This is a finite-horizon regularization effect, not new
address capacity and not an unconditional advantage of ill-conditioning.

\subsection{Branch reliability shrinkage}

Partition the parameter coordinates into disjoint branch blocks $b$. Let
$\bm g_b$ be the clean gradient in block $b$, $\Sigma_b$ its noise
covariance, and $I_b$ the identity on that block. Define signal and noise
energies $S_b=\lVert\bm g_b\rVert_2^2$ and
$N_b=\operatorname{tr}(\Sigma_b)$. With attenuation $a_b$ in each block,
$M=\operatorname{blockdiag}(a_bI_b)$ and a fixed step $\eta=c/L$ for $c>0$,
Eq.~\ref{eq:s_operator_bound} guarantees a loss decrease of
\begin{equation}
\mathcal G_c(\bm a)=\frac{c}{L}\sum_b
\left[a_bS_b-\frac{c a_b^2}{2}(S_b+N_b)\right].
\label{eq:s_branch_guarantee}
\end{equation}

\begin{proposition}[Fixed-step reliability-optimal attenuation]
For each block with $S_b+N_b>0$, Eq.~\ref{eq:s_branch_guarantee} is
maximized independently at
\begin{equation}
a_b^*=\min\left\{1,\frac{S_b}{c(S_b+N_b)}\right\},
\label{eq:s_reliability_gain}
\end{equation}
for a physical attenuator constrained to $0\leq a_b\leq1$. A block with
$S_b=N_b=0$ contributes zero for every gain and can be omitted from the
ratios below. If at least one block has positive energy, the best common gain is
$a_{\mathrm g}^*=\min\{1,\sum_bS_b/[c\sum_b(S_b+N_b)]\}$, and the
branch-specific guarantee is no smaller because its feasible set contains
every common gain.  At $c=1$, Eq.~\ref{eq:s_reliability_gain} reduces to the
reliability factor $S_b/(S_b+N_b)$; in that case
$\mathcal G_{\mathrm{branch}}=(2L)^{-1}\sum_bS_b^2/(S_b+N_b)$ and
$\mathcal G_{\mathrm{global}}=(2L)^{-1}(\sum_bS_b)^2/
\sum_b(S_b+N_b)$, with equality when reliability is constant across
positive-energy blocks. If all blocks have zero energy, every common gain
is optimal and both guarantees are zero.
\end{proposition}

\noindent\emph{Proof.}
Differentiate each concave quadratic in Eq.~\ref{eq:s_branch_guarantee} and
project its unconstrained maximizer onto $[0,1]$.  Restricting all coordinates
to a common value gives $a_{\mathrm g}^*$.  The branchwise optimum cannot be
worse than this restricted solution.  For $c=1$, the closed-form comparison
also follows from Cauchy--Schwarz applied to
$S_b/\sqrt{S_b+N_b}$ and $\sqrt{S_b+N_b}$. \hfill$\square$

This optimum assumes access to $S_b$ and $N_b$. A point implementation
provided with the same branch partition and multiplicative factors produces
the same update. Its biological relevance depends on whether local conductance can
estimate or realize those factors.

\subsection{Why depth can help or hurt}

Let $P_1\preceq P_2\preceq\cdots\preceq P_{D_{\mathrm{r}}}$ be nested tree-route
projectors. Increasing depth reduces the address bias
$\lVert(I-P_{D_{\mathrm{r}}})\bm g\rVert^2$ but increases admitted noise
$\operatorname{tr}(P_{D_{\mathrm{r}}}\Sigma P_{D_{\mathrm{r}}})$ whenever the new modes carry stochastic
energy. These competing terms can therefore produce an interior optimum
when task signal is concentrated above one hierarchy level and noise
continues into finer levels; its location depends on their magnitudes, as
specified by Eq.~\ref{eq:s_projection_crossover}. This is the controlled role isolated by
the hierarchy-depth routing experiment. It is distinct from nonlinear forward
composition, which can make deeper dendritic subunits representationally
useful even when the backward address span is unchanged.

Depth can also impose a transport cost. In an idealized ensemble with equal
local eligibility, let the route gain be
$A=\prod_{\ell=1}^{D_{\mathrm{p}}}G_\ell$ and assume independent Gaussian
log gains, $\log G_\ell\sim\mathcal N(\mu_\ell,\sigma_\ell^2)$. The cosine between a shared
coefficient and the exact gain-weighted field in the large-route limit is
\begin{equation}
\frac{\E A}{\sqrt{\E A^2}}
=\exp\left(-\frac12\sum_{\ell=1}^{D_{\mathrm{p}}}\sigma_\ell^2\right).
\label{eq:s_depth_transport_penalty}
\end{equation}
Additional levels can therefore add address resolution while simultaneously
making uncompensated shared feedback less aligned. The balance depends on
task hierarchy, gain variability, noise and the available coefficient
estimator; ``more depth'' has no unconditional sign.

\subsection{Controlled signal--noise experiment}

The signal--noise experiment used 16 Euclidean coordinates represented in the
orthonormal Haar basis of a balanced binary tree, ordered from the global mode
(level 0) to individual sibling contrasts (level 4). The configuration,
50 analysis seeds 7000--7049, two preliminary validation seeds and numerical
tolerances were specified before the analysis outcomes were inspected. The
preliminary runs verified basis orthogonality, projector idempotence,
static-gain span invariance and exact agreement with the explicit point-gain
control. All planned observations were completed: 3,750 spectral, 1,800
depth, 3,750 projection and 1,500 reliability rows.

In the spectral screen, the task-gradient second moment assigned total energy
$(0.15,0.35,0.30,0.20,0)$ to tree levels 0--4. An alignment coefficient
$\rho\in\{0,0.25,0.5,0.75,1\}$ mixed this second moment with an independently
rotated endpoint having the same eigenvalues. Intermediate mixtures do not
preserve the spectrum and jointly vary alignment and concentration.
Ancestry, random rank-matched
and dense principal-eigenspace projectors were compared at
$K\in\{1,2,4,8,16\}$. The dense value was also calculated as the Ky--Fan sum
of the leading $K$ eigenvalues divided by total second-moment energy. Because the
endpoint second moments had equal trace, ancestry-minus-random capture varied
affinely with alignment. Endpoint differences therefore gave the exact
per-seed crossover; seeds already favoring ancestry at zero alignment were
assigned threshold zero. In these unweighted coordinates, let
\(\overline\Sigma_g\) denote the uncentered task-gradient second moment; energy-weighted capture was
$\operatorname{tr}(P\overline\Sigma_g)/\operatorname{tr}(\overline\Sigma_g)$.

The same-span diagnostic compared raw nested indicators, their tree-Haar basis
and indicators multiplied by fixed nonzero lognormal column gains. Rank,
coherence, positive-spectrum Gram condition number and capture were calculated
directly.

In the depth screen, a unit-norm target was sampled from modes at levels
$0$ through $H_{\mathrm c}$, with $H_{\mathrm c}\in\{1,2,3,4\}$. The
gradient-noise standard deviation at level $\ell$ was
$0.65[1+1.5\max(\ell-H_{\mathrm c},0)]$. Starting from zero, 80
common-noise updates of size 0.12 used the projector onto levels $0$ through
$D_{\mathrm r}\in\{1,2,3,4\}$. A fixed leaf permutation provided the rewired
control, and the identity provided stochastic backpropagation.

The projection screen independently set retained signal $r$ and retained noise
$n$ to $\{0.25,0.5,0.75,0.9,1\}$. It compared an unprojected stochastic
gradient, backpropagation followed by the same projection, and a routed
estimator with an added total coefficient-noise fraction of 0.05. For a unit
step on the identity-curvature quadratic, expected losses were calculated
analytically and stochastic losses were recorded from paired draws.

In the reliability screen, eight disjoint blocks had signal and noise energies
proportional to $\exp(hx_b)$ and $\exp(-hx_b)$, respectively, with equally
spaced $x_b\in[-1,1]$ and heterogeneity
$h\in\{0,0.5,1,1.5,2\}$. We compared no attenuation, the best common gain,
$S_b/(S_b+N_b)$, seed-shuffled gains, reversed gains and a point-gain control receiving
the same branch factors. Intervals used 20,000 resamples of the 50 paired
seeds, and paired tests were two-sided Wilcoxon signed-rank tests.

The resolved protocol is
\texttt{configs/\allowbreak credit\_phase\_theory/\allowbreak confirmatory.json}.
\par\noindent Source tables are in
\texttt{source\_data/\allowbreak credit\_phase\_theory/}.
All full-rank projectors agreed to $4.44\times10^{-16}$, route-wise
static gain rescaling preserved the projector to $1.33\times10^{-15}$, and the
explicit point-gain control matched reliability-aligned attenuation exactly. This
fixed-state quadratic experiment validates the signal--noise prediction but does not join
oracle reliability to learned positive conductances.

\subsection{Initialization score in the trained subtree factorial}

The existing factorial used full-batch gradient descent. Its secondary
stochastic diagnostic reconstructed the frozen training data, initialization
and routes, then estimated routed-update variance from 64 hypothetical
minibatches of 64 examples per seed. The full-data routed mean and that
variance supplied Eq.~\ref{eq:s_operator_utility}. This variance was not
noise experienced by the full-batch optimizer and cannot by itself explain
the trained contrasts. The original released score used the largest logistic
Hessian eigenvalue at initialization as its curvature scale. The corrected
analysis uses the global logistic bound
\[
L_{\rm glob}=\max_c\lambda_{\max}\!\left(
\frac{8}{4N}X_c^{\mathsf T}X_c\right),
\]
where rows of $X_c$ contain the selected-stream features for examples in
context $c$, $N$ is
the total training-set size, and eight is the gradient context-rescaling
factor. This bound is 1.00083--1.02125 times the initialization Hessian
maximum across seeds. We scored the deterministic full-batch update with
$V=0$, the hypothetical minibatch estimator with its estimated $V$, and the
deterministic bound at the actual training step $\eta=0.3$.

Across 540 tree fits, deterministic utility correlated with norm-matched one-step progress at 0.9509 and final accuracy at 0.9242; hypothetical-minibatch values were 0.9373 and 0.9157. Intervals use 5,000 whole-seed resamples. These are retrospective associations within the factorial, not independent prospective validation.

To characterize the task-coordinate spectrum, the analysis used the
uncentered second moment of $\delta_i\bm e_{c_i}$, where $\delta_i$ is
the output error on example $i$ and $\bm e_{c_i}$ is the coordinate vector
for its context $c_i$. Eight balanced contexts and logits near zero make this spectrum
nearly isotropic, giving the reported participation rank near eight. This
value characterizes the chosen coefficient representation at initialization;
it is not an independent estimate of task complexity or a morphology law.
The mean within-family utility and accuracy maxima occur at $K=8$ for all
five nondegenerate route families. Across their 25 implementation--family
cells, there are nine distinct nonconstant mean utility curves: four
implementations coincide, whereas leaf rewiring changes four of the five
family curves; the principal component analysis (PCA) oracle curve is unchanged. In the dendritic-tree implementation,
deterministic utility selected an accuracy-maximizing bandwidth in 95 of 100
seed--family curves, the same success count as always choosing the largest
$K$; the hypothetical-minibatch score matched in 97 of 100. The 20
derangement curves had zero optimized-step utility throughout and were
reported as degenerate. Repeated equivalent parameterizations do not
provide independent tests of endpoint ordering. The distinct
interior-optimum and ancestry-minus-control analyses, including their
prediction failures, are reported in Section~\ref{note:interior_optimum}.
The corrected scores, per-family selection errors and regret are provided in
\texttt{source\_data/review\_evidence\_reanalysis/}.


\subsection{State-matched positive-conductance reliability experiment}

We next tested whether positive conductances could realize the attenuation
prescribed by Eq.~\ref{eq:s_reliability_gain} in a separate trained branch
model (Table~\ref{tab:source_index}). Here $b$ indexes branches, $i$ examples and $d$ inputs within a branch. Every
input rate $x_{bid}$ was nonnegative, and each excitatory conductance was
parameterized by $g_{bd}=\operatorname{softplus}(q_{bd})$, with trainable
raw parameter $q_{bd}$. With leak conductance $g^L$, define total unshunted
conductance and voltage on example $i$
\begin{equation}
G_{bi}=g^L+\sum_d x_{bid}g_{bd},\qquad
V_{bi}=\frac{\sum_d x_{bid}g_{bd}}{G_{bi}}.
\label{eq:s_positive_branch_voltage}
\end{equation}
A fixed shunt $\kappa_b\geq0$ had reversal zero. To isolate its effect on
input resistance from its effect on the forward feature, an oracle current
$\kappa_bV_{bi}$ was added on each example. The perturbed state is therefore
\begin{equation}
V_{bi}'=\frac{\sum_d x_{bid}g_{bd}+\kappa_bV_{bi}}{G_{bi}+\kappa_b}=V_{bi},
\qquad
\frac{\partial V_{bi}'}{\partial q_{bd}}
=\frac{x_{bid}(1-V_{bi})}{G_{bi}+\kappa_b}\,
\operatorname{sigmoid}(q_{bd}).
\label{eq:s_state_matched_shunt}
\end{equation}
Relative to the unshunted eligibility, the example-specific physical
input-resistance factor is $G_{bi}/(G_{bi}+\kappa_b)$. The compensating current was recomputed from the
pre-shunt state and treated as a clamp, not differentiated as an independently
learned pathway. Consequently this manipulation deliberately removes forward
feature suppression while retaining conductance attenuation in the local
update.

The binary task used eight branches, four inputs per branch, 512 training
examples and 2,048 held-out examples. Input rates were generated from a
positive baseline plus a class-dependent branch pattern and clipped at zero.
Initial full-data gradients supplied
$S_b=\lVert\bm g_b\rVert_2^2$. Coefficient-noise variance was then scaled so
that
\begin{equation}
N_b=S_b\exp(-2h x_b),\qquad x_b\in[-1,1],
\end{equation}
where $x_b$ is the evenly spaced branch-order coordinate and
$h\in\{0,0.5,1,1.5,2\}$ controls reliability heterogeneity. The statistical reliability was
$r_b=S_b/(S_b+N_b)$, lower-bounded at 0.03 for numerical stability. Because
the common step was $\eta=c/L$ with $c=0.5$, the corrected fixed-step target
was $a_b^*=\min\{1,r_b/c\}$ and the initial shunt was
$\kappa_b=\overline G_b[(a_b^*)^{-1}-1]$, where $\overline G_b$ is the
training-example mean of $G_{bi}$ at initialization. Shunts remained fixed during the 40
updates; the excitatory conductance parameters were trained, and the same
coefficient-noise realizations were used across matched methods within seed.

Nine conditions separated the relevant alternatives: exact clean
backpropagation, noisy credit without a shunt, a state-matched
unattenuated-eligibility control, the step-consistent best single global shunt,
reliability-aligned, seed-shuffled and anti-aligned shunts, an explicit
point-gain control receiving the identical branch factors, and aligned
shunting with clean credit. The point-gain control is a resource
control, not a biological assertion: because it receives the same branch
partition and factors, its update must equal Eq.~\ref{eq:s_state_matched_shunt}.
The clean-backpropagation condition retains the unattenuated exact gradient and
is the optimization reference.

An initial implementation using $r_b$ despite $c=0.5$ was excluded; the corrected protocol preceded two excluded checks and fifty paired seeds 8100--8149 (2,250 outcomes). State mismatch was below $3\times10^{-16}$; independently computed point-gain and unattenuated controls agreed to $1.4\times10^{-14}$, and fixed-compensating-current finite differences to relative error $3.60\times10^{-8}$.

At heterogeneity two, aligned-minus-best-global one-step test-loss reduction was 0.001817 (95\% paired interval, 0.001618--0.002020; 49/50 positive). Aligned attenuation also improved immediate progress over unshunted learning. After forty updates it remained better than global, shuffled and reversed attenuation, but its loss difference from no shunt was 0.000425 (95\% interval, $-0.002456$ to 0.003071). Clean BP retained the lowest mean loss. Thus the state-clamped oracle factors improve relative placement without establishing sustained superiority to unshunted learning; compact endpoints and all seeds are indexed below.

\subsection{Adaptive local reliability estimation}

The preceding experiment supplied the reliability moments. We next asked
whether they could be estimated from noisy local observations, specifying
the adaptive-gain experiment before examining its outcomes
(Table~\ref{tab:source_index}). At the current parameter
state, branch $b$ received two conditionally independent noisy local gradient
observations, $\widehat{\bm g}_{1b}$ and $\widehat{\bm g}_{2b}$, on the same
examples. It formed
\begin{equation}
 \widehat S_b^{\mathrm{obs}}
 =\langle\widehat{\bm g}_{1b},\widehat{\bm g}_{2b}\rangle,
 \qquad
 \widehat N_b^{\mathrm{obs}}
 =\frac12\lVert\widehat{\bm g}_{1b}-\widehat{\bm g}_{2b}\rVert_2^2.
\label{eq:s_adaptive_reliability_moments}
\end{equation}
For independent zero-mean credit noise, these estimate the squared clean
branch gradient and the noise energy of one observation. Eight paired
observations initialized the moments; exponential averages with decay 0.9
were then updated before each of 40 training steps. Negative finite-sample
signal estimates and moments below $10^{-12}$ were clipped only when the gain
was constructed. The step-consistent gain followed
Eq.~\ref{eq:s_reliability_gain}, was bounded to $[0.03,1]$, and was converted
to a shunt using the branch's current mean total conductance.

The task and state clamp are unchanged. At heterogeneity 0/1/2, seven rules compare clean BP, noisy no-shunt, fixed initial oracle, adaptive local/global/permuted shunts and the independently computed adaptive point gain. Update and probe noise are paired. Two implementation-check seeds are excluded; fifty paired seeds 8200--8249 supply 1,050 outcomes. State, finite-difference and point-gain checks pass at the stated tolerances.

At heterogeneity two, adaptive local attenuation lowers final loss relative to global and shuffled estimates, yet exceeds no-shunt loss by 0.004250 (95\% interval, 0.002499--0.005975) and fixed-oracle loss by 0.003449 (0.002577--0.004305). Only 26/50 immediate local-versus-shuffled differences favored local, failing the predefined direction criterion. Final gain ordering correlates with the initial oracle at 0.756, with mean absolute error 0.205. The supplied paired probes, optimizer step fraction and exact state clamp remain necessary assumptions; this is not a fully autonomous inhibitory-plasticity mechanism.

\subsection{Same-span noisy coefficient-learning experiment}

To isolate coefficient learning from address capacity, we compared three
dictionaries with the same span in a rate-based factorial
(Table~\ref{tab:source_index}). The 16-leaf balanced-tree Haar basis supplied the common
rank-eight span through depth three. The tested dictionaries were the eight
orthonormal Haar modes, all 15 nested indicators through depth three, and the
same nested indicators scaled by positive gains evenly spaced on the log scale
from $\exp(-3)$ to $\exp(3)$. The dictionaries' projectors agreed to
$9.44\times10^{-16}$ and their positive-spectrum Gram condition numbers were
1, 15 and 388.52.

Each independent seed drew a unit-norm target in the common span. At every
update the estimator observed the target plus isotropic Gaussian noise with
standard deviation $2/\sqrt{n_{\mathrm{eff}}}$, where
$n_{\mathrm{eff}}\in\{4,16,64,256\}$. The same 80 observation vectors were
used for every matched parameterization. Vanilla coefficient descent used
$0.2/\lambda_{\max}(\Phi^{\mathsf T}\Phi)$; the oracle coordinate control used
$0.2(\Phi^{\mathsf T}\Phi)^+$. Population field loss was recorded after 1, 5, 20
and 80 updates. The target address residual was below
$1.45\times10^{-30}$ in every condition, so all reported loss is within-span
coefficient-estimation error.

The predictions were specified before two implementation-check seeds, which
were excluded from inference, and 50 paired seeds (9100--9149). The results contradicted the prediction that Haar coordinates would
always reduce finite-time error. At $n_{\mathrm{eff}}=4$, raw-minus-Haar
final loss was $-0.2774$ (95\% paired interval $-0.3260$ to $-0.2313$) and
scaled-minus-Haar was $-0.1906$ ($-0.2500$ to $-0.1335$): slower modes
filtered high estimator noise. At $n_{\mathrm{eff}}=256$, the effects reversed
to 0.02555 (0.02171--0.02952) and 0.17091 (0.14576--0.19686), positive in
50/50 seeds for both contrasts. Equation~\ref{eq:s_coefficient_finite_risk},
derived after these outcomes to explain the reversal, predicted median
crossovers at effective
sample sizes 38.53 and 8.25. The Gram-preconditioned field trajectories agreed
across dictionaries to $1.15\times10^{-15}$. Thus static parameterization can
alter finite-data learning without altering the attainable credit field.
